\documentclass[11pt,a4paper]{article}
\usepackage[T1]{fontenc}
\usepackage[utf8]{inputenc}
\usepackage[danish]{babel}
\usepackage{lmodern,amsmath,amssymb,booktabs,array,microtype}
\usepackage[margin=25mm,headheight=15pt]{geometry}
\usepackage{xcolor,fancyhdr,hyperref}
\definecolor{ink}{HTML}{173F4F}
\hypersetup{colorlinks=true,linkcolor=ink,urlcolor=ink,pdftitle={Matematikken bag prismotor v4}}
\pagestyle{fancy}\fancyhf{}\fancyhead[L]{\small\textcolor{ink}{PRISMOTOR V4 / MATEMATISK NOTAT}}\fancyfoot[R]{\thepage}
\setlength{\parindent}{0pt}\setlength{\parskip}{6pt}
\newcommand{\E}{\mathbb E}\newcommand{\Prob}{\mathbb P}
\newcommand{\clip}{\operatorname{clip}}\newcommand{\pos}[1]{\left[#1\right]_+}
\newcommand{\code}[1]{\texttt{\small\detokenize{#1}}}
\begin{document}
{\Huge\bfseries\textcolor{ink}{Matematikken bag v4}}\\[7pt]
{\large Efterspørgselsfordeling, prisvalg og værdien af kapacitet}\\[5pt]
{\small Dokumenteret ud fra projektets reference og implementering, 3. oktober 2026.}

\section{Hvad optimerer modellen?}
Version 4 vurderer hvert muligt pristrin ud fra det forventede bidrag fra
\emph{fremtidigt} salg. Den vælger ikke en pris for at ramme en bestemt
målbelægning. For et lager og en opholdsdato er kerneproblemet
\begin{equation}\boxed{\quad k^*=\arg\max_k\;\underbrace{\bigl((1-c)p_k-v\bigr)}_{\text{bidrag pr. solgt enhed}}
\underbrace{\E[\min(X_k,C)]}_{\text{forventet fremtidigt salg}}.\quad}\label{eq:objective}\end{equation}
Her er $X_k$ den ikke-negative, heltallige efterspørgsel ved pris $p_k$, og $C$
er ledig kapacitet. Allerede bogført omsætning ændres ikke ved at vælge en ny pris
og indgår derfor ikke i denne optimering.

Koden kalder målet \code{expected_revenue}. Fordi både provision og variabel
omkostning trækkes fra, er det mere præcist et \textbf{forventet dækningsbidrag}
end bruttoomsætning. Der beregnes særskilt for private værelser og dormsenge.

\begin{center}\small
\begin{tabular}{@{}p{24mm}p{119mm}@{}}\toprule
Symbol & Betydning \\\midrule
$t,\ell$ & Opholdsdato og antal dage til opholdsdatoen (lead time).\\
$B,C$ & Enheder på bøgerne (OTB) og resterende salgsbar kapacitet.\\
$Y_t,R_{t,\ell}$ & Historisk endeligt salg og netto pickup: $R_{t,\ell}=Y_t-B_{t,\ell}$.\\
$p_k,p_0$ & Pris på trin $k$ og pris på referencetrinnet.\\
$c,v$ & Provisionsandel og variabel omkostning pr. solgt enhed.\\
$\varepsilon,L$ & Antaget priselasticitet og adaptivt efterspørgselsniveau.\\
$\mu,Z$ & Prognoseskala for pickup og empirisk formvariabel.\\
$S_D(x)$ & Halesandsynlighed $\Prob(D\geq x)$ for resterende efterspørgsel.\\\bottomrule
\end{tabular}
\end{center}

\textbf{Tre lag i beslutningen.} Først estimeres efterspørgslen ved
referenceprisen. Derefter beregnes økonomien på alle pristrin. Til sidst
anvendes markedsjustering, events, træghed og prisgrænser. Den endelige pris
behøver derfor ikke være det ubetingede optimum i (\ref{eq:objective}).

Dette er en model af de resterende salg ved en given pris, ikke en fuld dynamisk
optimering af alle fremtidige prisændringer, bookingforløb og ophold på én gang.

\newpage
\section{Fra reservationshistorik til efterspørgsel}
Bookingkuben rekonstruerer $B_{t,\ell}$ ved hjælp af oprettelses- og
annulleringsdatoer. Netto pickup kan være negativ, når annulleringerne overstiger
nysalget. Flytninger og andre ændringer uden fuld ændringshistorik kan ikke
rekonstrueres eksakt.

\subsection{Sæson og gennemsnit}
Lad $m(t)$ være måned og $w(t)$ ugedag. På de brugbare træningsdatoer estimeres
\begin{align}
s_m&=\frac{\operatorname{gennemsnit}(Y_t\mid m(t)=m)}{\operatorname{gennemsnit}(Y_t)},\\
a_{\ell,w}&=\operatorname{gennemsnit}\left(\frac{R_{t,\ell}}{s_{m(t)}}\;\middle|\;w(t)=w\right),\\
\mu^{\rm raw}_{t,\ell}&=\max\{0,a_{\ell,w(t)}s_{m(t)}L\}.
\end{align}
Sæsonfaktoren er dimensionsløs. $a_{\ell,w}$ og $\mu^{\rm raw}$ måles i enheder.
Udsolgte datoer flagges som censurerede og udelades som standard; særlige
lukkeperioder kan også udelades. Det fjerner ikke censureringsproblemet: de
udeladte datoer kan netop være dem med størst latent efterspørgsel.

\subsection{Den empiriske form}
For observationer med $a_{\ell,w}\geq2$ beregnes forholdet
\begin{equation}z_{t,\ell}=\frac{R_{t,\ell}}{s_{m(t)}a_{\ell,w(t)}}.\end{equation}
Disse forhold samles i lead-intervallerne
\[
[0,1],\ [2,3],\ [4,6],\ [7,10],\ [11,17],\ [18,30],\ [31,45],\ [46,70],\ [71,365].
\]
Hvert interval gemmes som 201 empiriske kvantiler $q_0,\ldots,q_{200}$.
Det bevarer skævhed, gruppespring og negative observationer uden at vælge en
Poisson-fordeling. Tomme intervaller låner en form fra et andet interval.

Mellem to forskellige nabokvantiler bruges lineær interpolation. Hvis
$q_{i-1}<z\leq q_i$, er den implementerede hale
\begin{equation}
S_Z(z)=1-\frac{i-1+(z-q_{i-1})/(q_i-q_{i-1})}{200}.
\end{equation}
Gentagne kvantiler håndteres særskilt i koden. Over højeste kvantil bruges en
konstrueret hale, med $T=\max(q_{200},10^{-9})$:
\begin{equation}
S_Z(z)=\max\left\{10^{-4},\frac{1}{201}\exp\left(-\frac{z-T}{\max(0{,}25T,10^{-9})}\right)\right\}.
\end{equation}
Dette er en numerisk regularisering, ikke en estimeret ekstremhale. Gulvet og
springet lige over største kvantil betyder, at funktionen ikke overalt er en
matematisk konsistent overlevelsesfunktion. Forventningsformlerne nedenfor er
idealet; koden bruger denne numeriske approksimation.

\newpage
\section{Krympning mod kalendergennemsnittet}
En stor tidlig gruppebooking er ikke nødvendigvis et pålideligt signal om det
endelige salg. Derfor kombineres bøger plus pickup med et kalendergennemsnit
$H_{w,m}=\operatorname{gennemsnit}(Y_t\mid w,m)$:
\begin{align}
\widehat Y_{t,\ell}&=w_\ell(B_{t,\ell}+\mu^{\rm raw}_{t,\ell})
 +(1-w_\ell)H_{w(t),m(t)}L,\label{eq:shrink}\\
\mu_{t,\ell}&=\max\{0,\widehat Y_{t,\ell}-B_{t,\ell}\}.\label{eq:mu}
\end{align}
$w_\ell=1$ lægger al vægt på bøger plus pickup; $w_\ell=0$ lægger al vægt på
kalenderen. Bindingen \code{model.bind(otb)} gør det aktuelle $B$ tilgængeligt
for efterspørgselsberegningen.

\subsection{Udledning af vægten}
På træningsdata sættes $A_i=B_i+\mu_i^{\rm raw}$ og $H_i=H_{w(i),m(i)}$ med
niveau $L=1$. For hvert lead minimeres
\begin{equation}\min_{0\leq w\leq1}\sum_i\bigl(H_i+w(A_i-H_i)-Y_i\bigr)^2.\end{equation}
Differentiering med hensyn til $w$ giver
\begin{equation}\boxed{w_\ell=\clip_{[0,1]}
\left(\frac{\sum_i(A_i-H_i)(Y_i-H_i)}{\sum_i(A_i-H_i)^2}\right).}\end{equation}
Hvis nævneren er for lille, gemmes ingen vægt for det lead. Koden bruger den
nærmeste gemte vægt eller $1$ uden krympningsmodel.

\subsection{Fordelingen efter krympning}
Den krympede skala multipliceres med den samme empiriske form:
\begin{equation}
D=\mu_{t,\ell}Z,\qquad S_D(x)=S_Z(x/\mu_{t,\ell}),\quad x>0,\ \mu_{t,\ell}>0.
\end{equation}
Ved $\mu=0$ sættes halen til nul for positive tærskler. Negative formkvantiler
kan stadig forekomme, men prisberegningen tæller kun positive salgsgrænser.

\textbf{Præcis fortolkning.} Krympningen skalerer fordelingen; den blander ikke
to fulde prognosefordelinger. $\mu$ er heller ikke nødvendigvis præcis
$\E[D]$: den poolede kvantilform er ikke tvunget til at have middelværdi $1$.
Desuden tillader (\ref{eq:mu}) ikke negativ forventet resterende pickup, selv om
træningsdata kan indeholde negativ pickup. Modellen er derfor ikke en fuld
stokastisk model for annulleringer og frigivelse af allerede solgt kapacitet.

\textbf{Eksempel.} Med $B=20$, $\mu^{\rm raw}=12$, $H=24$, $L=1$ og
$w=0{,}4$ bliver $\widehat Y=0{,}4\cdot32+0{,}6\cdot24=27{,}2$ og
$\mu=7{,}2$. Kalenderen dæmper dermed pickup-skalaen fra $12$ til $7{,}2$.

\newpage
\section{Priselasticitet og forventet salg}
Ved pris $p_k$ skaleres efterspørgslen med
\begin{equation}
g_k=\left(\frac{p_k}{p_0}\right)^{-\varepsilon},\qquad
D_k=g_kD,\qquad S_{D_k}(x)=S_D(x/g_k).
\end{equation}
De konfigurerede elasticiteter er $1{,}6$ for rum og $2{,}0$ for senge. De er
\emph{antagelser}, ikke målte kausale effekter på egne gæster. Eksempelvis giver
10\,\% højere værelsespris $g=1{,}1^{-1{,}6}\approx0{,}859$.

\subsection{Hvorfor summeres halerne?}
For ikke-negativ heltallig efterspørgsel $X$ og heltallig kapacitet $C$ gælder
\begin{align}
\min(X,C)&=\sum_{j=1}^{C}\mathbf1\{X\geq j\},\\
\E[\min(X,C)]&=\sum_{j=1}^{C}\Prob(X\geq j).
\end{align}
Koden evaluerer derfor
\begin{equation}
Q_k(C)=\sum_{j=1}^{C}S_D(j/g_k),\qquad V_k(C)=\bigl((1-c)p_k-v\bigr)Q_k(C).
\end{equation}
For en kontinuert $D_k$ svarer summen til at tælle hele positive enheder, dvs.
$X_k=\lfloor\max(D_k,0)\rfloor$, hvis halen var en konsistent fordeling.
Den kontinuerte forventning ville i stedet være integralet
$\int_0^C S_{D_k}(x)\,dx$.
\code{expected_sales} afrunder kapaciteten op med en numerisk tolerance og
begrænser som standard summen til 400 enheder.

\subsection{Et gennemregnet prisvalg}
Et rent regneeksempel: $C=3$, $c=0{,}16$, $v=95$ kr. Tabellen antager de viste
halesandsynligheder ved hvert trin; de er ikke en prognose for hostellet.
\begin{center}\small
\begin{tabular}{@{}rrrrrrr@{}}\toprule
$p_k$ & $S_k(1)$ & $S_k(2)$ & $S_k(3)$ & $Q_k$ & $(1-c)p_k-v$ & $V_k$\\\midrule
600 & 0,95 & 0,80 & 0,55 & 2,30 & 409 & 940,70\\
700 & 0,90 & 0,70 & 0,40 & 2,00 & 493 & 986,00\\
800 & 0,80 & 0,55 & 0,25 & 1,60 & 577 & 923,20\\\bottomrule
\end{tabular}
\end{center}
700 kr. vinder: højere bidrag pr. salg opvejer tabet af volumen fra 600 kr.,
men ikke det yderligere volumetab ved 800 kr.

Uden kapacitetsbinding giver den kontinuerte approksimation
$V(p)\propto((1-c)p-v)p^{-\varepsilon}$. For $\varepsilon>1$ har den optimum
\begin{equation}p^*=\frac{\varepsilon v}{(\varepsilon-1)(1-c)}.\end{equation}
Med rumparametrene fås ca. 302 kr. Det illustrerer, hvorfor en antaget høj
elasticitet kan presse priser ned: indtil måling begrænses det rå trinvalg til
højst ét trin under reference. Dette er en politikbegrænsning, ikke et estimat.

\newpage
\section{Bid price: værdien af den sidste enhed}
Ved fast pris og bidrag $n=(1-c)p-v\geq0$ er værdien af kapacitet $C$
$V(C)=n\sum_{j=1}^{C}S_D(j/g)$. Den marginale værdi er derfor
\begin{equation}\boxed{b(C)=V(C)-V(C-1)=nS_D(C/g).}\end{equation}
Bid price er den forventede alternativomkostning ved at afgive den sidste
ledige enhed. I eksemplet med 700 kr. er $b(3)=493\cdot0{,}40=197{,}20$ kr.
Koden bruger $\max(0,n)$ og returnerer dette bidrag ved $C\leq0$ som en
konvention; den konvention er ikke en tilladelse til overbooking.

\subsection{Flex-rum}
Lad $f$ være antal flex-rum allokeret til private, og lad et flex-rum rumme $h$
senge. Kapaciteterne bliver $C_r(f)=C_r(0)+f$ og $C_b(f)=C_b(0)-hf$.
Ved faste priser og efterspørgselsmodeller er gevinsten ved ét ekstra privat rum
\begin{align}
G(f)&=n_rS_r\left((C_r(f)+1)/g_r\right),\\
T(f)&=n_b\sum_{j=0}^{h-1}S_b\left((C_b(f)-j)/g_b\right),\\
\Delta(f)&=G(f)-T(f).
\end{align}
Flytningen foretages, når $\Delta(f)>\tau$, hvor $\tau=25$ kr. i
konfigurationen. Her er $h=8$; det tidligere lagerekssempel i referencen bruger
4 og er ikke den aktuelle parameter.

Algoritmen starter ved den tilladte nedre grænse og lægger rum til, indtil
gevinsten er for lille eller øvre grænse nås. For monotone haler falder
rumgevinsten, mens tabet på de fjernede senge stiger. Derfor er den grådige
regel begrundet for dette separable problem. Den numeriske ekstremhale omtalt
på side 2 betyder, at monotonicitet ikke er en ubetinget garanti i koden.

Trinvalget beregnes i motoren med hvert lagers \emph{maksimalt mulige}
kapacitet. Begge lagre kan ikke samtidig få alle flex-rum. Prisvalg og
flex-allokering er derfor ikke én fælles global optimering. Rumtyper,
faktiske reservationer og ophold på tværs af datoer begrænser desuden
allokeringen.

\subsection{Ophold over flere nætter}
For ét ophold over datoerne $\mathcal T$ er den implementerede kapacitetsværdi
\begin{equation}B_{\rm ophold}=\sum_{t\in\mathcal T}b_t(C_t).\end{equation}
Et tilbuds samlede bidrag kan sammenlignes med denne sum. Funktionen
\code{stay_value} findes, men er ikke koblet til bookingmotoren. Summen alene
løser ikke sammenhængende rumtildeling eller automatisk minimumsophold.

\newpage
\section{Grupper: forventet fortrængning og punktprognose}
Hvis en gruppe bruger $q\leq C$ enheder på en dato, er den marginale
fortrængningsværdi ved fast pris
\begin{equation}
A(q)=V(C)-V(C-q)=\sum_{i=0}^{q-1}b(C-i).
\end{equation}
De første enheder kan koste lidt, mens de sidste koster meget. Med eksempellets
700 kr. koster to enheder $493(0{,}40+0{,}70)=542{,}30$ kr. i mistet
forventet dækningsbidrag.

\textbf{Motorens konkrete gruppegulv er mere konservativt.} I
\code{engine.py} anvendes $\widetilde n_t=(1-c)p_t$ i fortrængningssummen:
der fratrækkes ikke variable omkostninger på det fortrængte salg. Derefter
lægges gruppens variable omkostning til:
\begin{align}
\widetilde A_t(q)&=(1-c)p_t\sum_{i=0}^{q-1}S_t\left((C_t-i)/g_t\right),\\
r_{t,\min}&=v+\frac{\widetilde A_t(q)}q.\label{eq:group}
\end{align}
Med to rum i eksemplet giver dette $v+588\cdot1{,}10/2=418{,}40$ kr. pr. rum,
før prisgrænser og afrunding. Et gulv baseret på mistet dækningsbidrag ville
i stedet være $95+542{,}30/2=366{,}15$ kr.

Som standard sættes $g_t=1$ i gruppeberegningen, så længe elasticiteten ikke er
målt. Den særskilte hjælpefunktion \code{group_floor} kan gennemsnitsberegne
over flere nætter med de nettopriser, den får; motorens viste tilbud beregnes
pr. dato med (\ref{eq:group}). Der indgår ingen særskilt division med
$1-c_{\rm gruppe}$ for en eventuel provision på gruppetilbuddet.

\subsection{Det supplerende punktbaserede tal}
Motoren viser også \code{minimum_rate_certain}. Med total rumkapacitet $K_t$
og punktprognose $\widehat F_t$ er beregningen
\begin{equation}
r_{t,\rm punkt}=v+\frac{(1-c)p_t}{q}
\max\{0,q-\max(0,K_t-\widehat F_t)\}.
\end{equation}
Navnet ``certain'' betyder her en beregning, der tager punktprognosen som
udfald. Det er hverken en øvre fraktil, et worst case eller en garanti.
Begge tal underlægges prisgulv, prisloft og afrunding.

\section{Adaptivt efterspørgselsniveau}
Forventet og faktisk pickup summeres først på tværs af datoer til $E$ og $A$.
Den implementerede opdatering, når $E\geq20$, er
\begin{equation}\boxed{L_{\rm ny}=\clip_{[0{,}6,1{,}6]}
\left[L_{\rm gammel}\left(\frac{\max(0,A)+1}{E+1}\right)^{0{,}07}\right].}\end{equation}
Ved $E<20$ ændres niveauet ikke. Tillægget $1$ undgår nulstilling ved nul
pickup. For $L=1$, $E=100$ og $A=75$ fås ca. $0{,}980$: en lille korrektion.
Rum og senge har separate niveauer. Dette er en adaptiv regel, ikke en ny
estimation af elasticiteten eller hele fordelingen.

\newpage
\section{Fra økonomisk optimum til endelig pris}
De aktuelle trinfaktorer er
\[
(0{,}80;\ 0{,}86;\ 0{,}93;\ 1{,}00;\ 1{,}07;\ 1{,}15;\ 1{,}24;\ 1{,}35;\ 1{,}48).
\]
De multipliceres med dagens grundpris og behandles med prisgrænser og
afrunding. Ved ens målværdi vælger \code{choose_rung} det første trin.

Efter den indledende rabatbegrænsning justeres trinpositionen med markedet:
\begin{equation}
u=k_{\text{begrænset}}+\frac{w_M}{s}\clip_{[-1,1]}
\left(\frac{p_{\rm comp}I_{\rm kvalitet}/p_{\rm base}-1}{a_M}\right).
\end{equation}
Her er $w_M=0{,}20$, $s=0{,}15$ og $a_M=0{,}20$. Uden markedspris er
markedsleddet nul. Mere end tre dage før ankomst udglattes positionen med
forrige position, når den findes: $\bar u=0{,}5u+0{,}5u_{\text{før}}$.
Derefter tilføjes eventtrin $\operatorname{round}(\max(0,F_{\rm event}-1)/0{,}10)$.

Hysterese kræver normalt en afstand på mindst $0{,}75$ trin før skift.
Der tillades højst to trin op og ét ned pr. trinbeslutning; tæt på ankomst to
ned. Desuden gælder bl.a. eventgulv, langtidsgrænse, manuelle låse, prisgulv,
prisloft og en daglig ændringsbremse på 15\,\% med fast 24-timers anker.
Det særskilte knaphedsgulv og ratchet er slået fra ved aktiv v4.

\textbf{Eksploration i koden.} Et deterministisk hashvalg pr. dato og lager kan
vælge et nabotrin ved lead højst 60 dage, hvis afstanden mellem positionen og
det valgte trin er højst $0{,}35$. Dette tester ikke direkte, om to trins
forventede dækningsbidrag er næsten ens. Referencens beskrivelse af ``lige
gode trin'' er derfor en hensigt, ikke en verificeret økonomisk betingelse.
En kausal elasticitetsmåling kræver også et korrekt forsøgsdesign og logning af
de priser, gæsterne faktisk blev eksponeret for.

\section{Hvad kan valideres?}
Prognosepræcision kan måles på senere, adskilte historiske perioder med fx.
$\mathrm{MAE}=N^{-1}\sum_t|Y_t-\widehat Y_t|$ og empirisk dækning af
prognoseintervaller. Gentagne lead-observationer fra samme dato er afhængige;
derfor bør usikkerhed beregnes med datoen som klynge. Referencens historiske
resultater er ikke genkørt som del af dette dokument.

En bedre salgsprognose dokumenterer ikke en omsætningsgevinst ved andre priser.
Historiske priser afhænger selv af forventet efterspørgsel. Modellen rummer
også andre antagelser end elasticiteten: censurering, sæsonstruktur,
formpooling, haleapproksimation og driftsregler påvirker beslutningen.

\vfill
{\small\raggedright\textbf{Kildegrundlag og reproduktion.}
Projektets \code{REFERENCE.md} samt \code{app/cube.py}, \code{app/demand.py},
\code{app/bidprice.py}, \code{app/engine.py}, \code{app/ladder.py},
\code{app/level.py} og \code{config/config.yaml}. Ved forskel mellem forklarende
tekst og kode beskriver notatet koden. Alle regneeksempler er illustrative.
LaTeX-kilden kompileres med \code{pdflatex v4-matematik.tex} (to kørsler).\par}
\end{document}
